\documentclass{beamer}
\usepackage[spanish]{babel}
\usepackage[utf8]{inputenc}
\usepackage{graphicx}

\newtheorem{procedimiento}{Cabe destacar}
\newtheorem{forma}{¿Como miramos la eficiencia de la interpolación de Taylor?}
\newtheorem{idea}{¿En que consiste la interpolación de Taylor?}
\newtheorem{descripcion de los experimentos}{\underline {Descripcion de los experimentos}}
\newtheorem{introduccion}{\underline {Introducción}}
\newtheorem{objetivos}{\underline {Objetivos}}
\newtheorem{material}{\underline {Material utilizado}}
\newtheorem{resultados}{\underline {Resultados}}
\newtheorem{analisis}{\underline {Analisis de los resultados}}
\newtheorem{anotacion}{ }

%%%%%%%%%%%%%%%%%%%%%%%%%%%%%%%%%%%%%%%%%%%%%%%%%%%%%%%%%%%%%%%%%%%%%%%%%%%%%%%
\title[Presentación con Beamer]{Trabajo final: Interpolación de Taylor.}
\author[Técnicas Experimentales]{Bolaños Florido,Cynthia. \\
Cruz Guerra, Adrián.\\
Díaz Rodríguez, Diego}
\date[15-05-2013]{Miércoles 15 de marzo de 2013}
%%%%%%%%%%%%%%%%%%%%%%%%%%%%%%%%%%%%%%%%%%%%%%%%%%%%%%%%%%%%%%%%%%%%%%%%%%%%%%%

%\usetheme{Madrid}
%\usetheme{Antibes}
%\usetheme{boxes}
\usetheme{tree}
%\usetheme{classic}

\begin{document}

\begin{frame}

\includegraphics[width=0.15\textwidth]{ullesc.eps}
\hspace*{7.5cm}
\includegraphics[width=0.16\textwidth]{fmatesc.eps}
\titlepage

  \begin{scriptsize}
    \begin{center}
     Facultad de Matemáticas \\
     Universidad de La Laguna \\
     Grupo 1I.
    \end{center}
  \end{scriptsize}
\end{frame}


\begin{frame}
  \frametitle{Índice}  
  \tableofcontents[pausesections]
\end{frame}


\section{Objetivos.}
\begin{frame}
\frametitle{Objetivos.}
 \tableofcontents[pausesections]
\begin{introduccion}
El siguiente informe científico-técnico ha sido realizado para exponer y explicar la interpolación de una función. 
En particular, dicha interpolación se realizará mediante el método de Taylor.
 \tableofcontents[pausesections]
\end{introduccion}
\begin{objetivos}
\begin {itemize}
\item \underline {Objetivo principal:} Interpolación  de la función específica\\ 
f(x)= 1/(4x) mediante el método de Taylor.\\
\item \underline {Objetivo secundario:} Poner en práctica los conocimientos adquiridos sobre Latex y en particular con beamer.
\end {itemize}
\end{objetivos}
\end{frame}


\section{Procedimiento experimental.}
\begin{frame}
\frametitle{Procedimiento experimental.}
\begin {descripcion de los experimentos}
 \tableofcontents[pausesections]
\begin{idea}
\begin {itemize}
\item La fórmula de Taylor consiste en aplicar la sucesión de Taylor:
Introducir la formula de Taylor. 
\end {itemize}
 \tableofcontents[pausesections]
\end {idea}
\begin{forma}
\begin {itemize}
\item Comparamos la función dada con la función interpolada.
Introducir la formula de Taylor y la función dada.
\end {itemize}
\end {forma}
\end{descripcion de los experimentos}
\end{frame}


\begin{frame}
\begin{procedimiento}
 \tableofcontents[pausesections]
\begin {itemize}
\item Hemos hallado la función interpolada según varios grados con el objetivo de analizar si ello influye en la exactitud de este método de aproximación.
Luego hemos hallado el error entre la función inicial y la función interpolada varias veces y además en varios puntos de x y c para saber si ello influye.
\end {itemize}
\end {procedimiento}
\end {frame}


\begin {frame}
\begin{material}
El material utilizado ha sido el siguiente:
 \tableofcontents[pausesections]
\begin {itemize}
\item COMPLETAR
\item COMPLETAR
\item COMPLETAR
\item COMPLETAR
\item COMPLETAR
\end {itemize}
\end {material}
\end{frame}


\section{Conclusiones.}
\begin{frame}
\frametitle{Conclusiones.}
 \tableofcontents[pausesections]
\begin{resultados}
Los resultados obtenidos se recogen en la siguiente tabla:
fórmula inicial  fórmula interpolada  error
\end{resultados}
\begin{analisis}
A la vista de la tabla anterior se concluye que:
\begin {itemize}
\item El polinomio de Taylor es más exacto cuanto mayor sea el grado de la derivada en la sucesión.
\item El valor de x influye?
\item EL valor de c influye?
\item Si c y x poseen el mismo valor entonces el error es 0.
Realizar el experimento 4 veces como min cambiando dichos valores para sacar conclusiones.
\end {itemize}
\end{analisis}
\end {frame}

\begin {frame}
\begin{procedimiento}
La interpolación de Taylor no es uno de los métodos más exactos de aproximación en comparación  a otros métodos como Lagrange.
\end {procedimiento}
\end{frame}

\section{Bibliografia}
\begin{frame}
 \frametitle{Bibliografía}
\begin {anotacion}  
   La información requerida para la realización de este informe ha sido extraída de las siguientes fuentes:
\end {anotacion}
  \begin{thebibliography}{10}
    \beamertemplatebookbibitems
    \bibitem[Bibliografía]{listado}
\begin{itemize}
    \item {\small $http://www.tat.physik.uni-tuebingen.de/~kokkotas/Teaching/Num_Methods_files/Comp_Phys3.pdf$}
    \item {\small $http://www.fceia.unr.edu.ar/lcc/cdrom/Instalaciones/LaTex/latex.html$}
\end {itemize}
  \end{thebibliography}
\end{frame}


\end{document}









